\FloatBarrier
\section{本章小结}
\label{sec:set-functions-chapter-summary}

本章从传感器覆盖出发，把对整组选择的评价写成有限底集上的集合函数。底集可以是传感器，也可以是图的顶点或边，不必带图结构。$n$元底集的完整取值表有$2^n$项，所以函数通常由求值规则给出，算法的代价要分开计入求值次数和每次求值的开销。指示向量把子集换成零一向量，只换了索引方式，不带来可加性或局部分解。边际增益$\Delta(v\mid S)$是后面一切判别的基本量：它依赖已选集合，沿任一加入次序累加，总和等于两端函数值之差。

用边际增益可以分别回答三个问题。单调性问新增元素会不会降低评价，等价于全部边际增益非负；模性问每个元素的贡献是否固定，等价于四值等式，模函数恰是常数加固定权重之和；次模性问收益是否随已选集合增大而递减，在有限底集上与四值不等式等价。覆盖量与非负边权的割权是次模的典型例子，协同效应则给出超模函数。

这些性质互不蕴含。割权次模而不单调；$\max\{3\,\mathbb I\{\{1,2\}\subseteq S\},\mathbb I\{3\in S\}\}$中元素$1,2$只有同时选中才有大价值，它单调而不次模；模函数同时次模和超模，却可能因为负权重而不单调。归一化只是减去常数$F(\varnothing)$的约定，不改变任何边际增益。所以后面引用某种算法保证时，要逐项写明它需要哪些性质，不能只凭“次模”判断。

子集累计与Möbius反演回答另一个问题：由全部累计值恢复每个子集自身的贡献。反演按元素个数之差交替加减，依靠的是组合数交错和为零的抵消，它与累计互为逆变换，结果可以为负。这一问题只依赖有限求和与包含关系，不需要单调或次模。计算上，逐个套用公式和按子集大小递推都要$3^n$次运算，逐元素差分降到$n2^{n-1}$次，但仍随底集大小指数增长。一般有限偏序集上，Möbius函数由区间递推确定，子集格上的交替符号是它的特例。

下一章利用本章的边际增益和次模性，在基数或预算约束下最大化集合函数；第~\ref{chap:pairwise-binary-minimization}章再把一类可分解为一元项与成对项的次模函数化为割，精确求最小值。
